\section{AI 社交网络 vs 互联网社交}

Elys 被问到是不是 AI 版 Tinder。Tristan 的回答是：它不是。Tinder 的核心是低维匹配和滑动选择，Elys 想做的是高维 context 匹配和分身主动社交。传统互联网社交需要人主动表达、筛选、破冰；AI 社交希望分身先替你表达和探索，再把可能的真实连接带回来。

\begin{figure}[H]
\centering
\includegraphics[width=0.92\textwidth]{ai-social-vs-internet-social.png}
\caption{互联网社交与 AI 社交网络的差异。}
\end{figure}

\begin{knowledgebox}{读图：AI 社交的差异不只是“多了 AI”}
左侧互联网社交依赖低维标签和用户主动行为；右侧 AI 社交依赖高维 context 和分身主动行动。真正差异在于连接成本被 AI 降低，而不是界面换成聊天框。
\end{knowledgebox}

\subsection{为什么不是 AI 版 Tinder}

Tinder 更像快速筛选关系候选，强调外显资料和即时选择。Elys 更想做 context 网络：一个人非常细微、甚至不适合公开说的价值观、癖好、表达方式，也可能通过 AI 找到共鸣者。它追求的是更低摩擦、更高维的连接。

\begin{warningbox}{AI 社交的危险误区}
如果 AI 社交只是生成虚假热闹，它会快速消耗信任。用户可以接受 AI 辅助破冰，但真正价值仍然来自真人和真人之间的连接。
\end{warningbox}

\subsection{谁拥有最多 Context}

访谈提到 ChatGPT、Meta、豆包等可能拥有大量 context。Tristan 的判断是，智能会逐渐平权，但 context 不平权。模型公司有模型能力，社交平台有关系链，ChatGPT 有个人记忆和对话历史，AI 社交创业公司则试图围绕 context 飞轮做专门产品。

\begin{figure}[H]
\centering
\includegraphics[width=0.92\textwidth]{context-competition.png}
\caption{AI 社交竞争变量：模型能力、社交关系链、Context 资产与用户主体性。}
\end{figure}

\begin{knowledgebox}{读图：为什么 Context 比模型更稀缺}
图中四类玩家各有优势：模型公司有智能，社交平台有关系，创业公司有新范式，用户本人有主体性。若智能逐渐商品化，长期 context 和真人关系会成为更稀缺的资产。
\end{knowledgebox}

\subsection{本章小结}

AI 社交网络的核心不是匹配算法换代，而是从低维标签社交转向高维 context 社交。它的目标是让 AI 降低真人连接成本，而不是替代真人连接。

